\section{A precise replacement for P14.8}
\label{sec:p148-precise-replacement}

Let \(X\) be a connected compact K\"ahler manifold, and let
\(\theta\) be a smooth closed real \((1,1)\)-form such that
\(\PSH(X,\theta)\neq\varnothing\).  Write
\[
 \mathscr D_{\mathcal I}(X,\theta)
 :=\PSH(X,\theta)/{\sim_{\mathcal I}}
\]
for the set of realizable \(\mathcal I\)-data.  Thus an element arises from
an actual \(\varphi\in\PSH(X,\theta)\) and is represented by the tower
\[
 \bigl(\mathcal I(m\varphi)\bigr)_{m\geq1},
\]
equivalently by its divisorial Lelong profile.  It is not an arbitrary
formal sequence of ideals.  For
\(d=[\varphi]_{\mathcal I}\), set
\begin{equation}
 \mathcal R_\theta(d):=P_\theta[\varphi]_{\mathcal I}.
 \label{eq:p148-reconstruction-map}
\end{equation}

\begin{theorem}
\label{thm:p148-reconstruction-dichotomy}
The following two assertions give the exact reconstruction dichotomy.

\begin{enumerate}
\item The map in \eqref{eq:p148-reconstruction-map} is well defined and
returns the unique normalized \(\mathcal I\)-model potential in the class
\(d\).  Thus full realizable \(\mathcal I\)-data reconstructs its canonical
\(\mathcal I\)-model representative, without any hypothesis on
nonpluripolar mass.

\item Raw nonpluripolar mass is not part of \(\mathcal I\)-data.  On
\(X=\mathbf P^1\), with the Fubini--Study form \(\omega\) normalized by
\(\int_X\omega=1\), one \(\mathcal I\)-class contains potentials
\(\varphi_s\), \(0\leq s\leq1\), such that
\begin{equation}
 [\varphi_s]_{\mathcal I}=[0]_{\mathcal I},
 \qquad
 \int_X\left\langle\omega+\ddc\varphi_s\right\rangle=1-s.
 \label{eq:p148-same-data-all-masses}
\end{equation}
In particular, the class of the raw zero-mass representative
\(\varphi_1\) has canonical normalized \(\mathcal I\)-model representative
\(0\), whose mass is one.  Hence neither a raw zero-mass representative nor
its mass can be reconstructed from its \(\mathcal I\)-class.
\end{enumerate}
\end{theorem}

The two mathematically defined readings of the former P14.8 wording are
therefore settled.  Full realizable \(\mathcal I\)-data determines the
canonical normalized \(\mathcal I\)-model representative, whereas neither a
raw representative nor its raw mass descends to the
\(\mathcal I\)-quotient.  Any different reconstruction problem must first
specify additional input and is not asserted here.

On a smooth projective variety, full realized DXZ--Boucksom--Jonsson objects
in a pseudoeffective real N\'eron--Severi class have their separate all-mass
order isomorphism; the literal all-semivaluation ideal formula is used only
for a \(\mathbf Q\)-line-bundle class.  No compact-K\"ahler or genuinely
transcendental extension is asserted here.  Realization of arbitrary formal
ideal or valuative data remains P14.1.  The terminal-left-limit corollary
isolated in P14.7 assumes a complete existing filtration curve and is not
used here.
